A coin is a small thing to place at the centre of a confrontation. It fits in a hand. It can pass from one person to another without either person looking closely at it. But Jesus asks his questioners to look. There is an image on this coin. There is an inscription. Whose are they?

They know the answer. Caesar's. What they have not understood is how much their answer will ask of them. They want Jesus to settle a question about money. His answer will reach the person holding it.

The question has been carefully prepared. These men praise Jesus for teaching the truth and for refusing to be swayed by appearances. Then they ask whether it is lawful to pay tribute to Caesar. A refusal could be used against him before the Roman authorities. An agreement could be used to turn the people against him. Their praise of truthfulness serves an untruthful purpose. They want an answer they can use, whichever answer he gives.

Jesus makes them identify the coin. Then he tells them to render what belongs to Caesar, and what belongs to God. Saint John Chrysostom draws attention to that word, render. Something is owed. Jesus has not made the actual obligation disappear. Belonging to God does not give us leave to treat the common life as somebody else's responsibility.

Yet Chrysostom also insists that an impious demand does not become right because Caesar makes it. Earthly authority has a real place. It does not have God's place. We must ask what is truly owed, and we must give it truthfully. We cannot let power decide for us what goodness means.

That may sound like a way of dividing life into portions: some money for Caesar, some time for God, and the rest for ourselves. But consider the question Jesus has opened. The coin bears Caesar's image. What about the person who carries the coin?

Saint Thomas Aquinas takes that further step in reading this Gospel. The soul bears God's image. The person is to be returned to God. The comparison changes the scale of the answer. Caesar can receive a coin. God calls for a life. And the life that belongs to God includes the hand that pays, the mind that judges, and the heart that decides why it is paying.

We cannot finish with God by assigning him one protected corner of the week. We belong to him when we work, when we speak about another person, when we carry a responsibility that nobody applauds. His claim reaches the way we fulfil our other duties. The question is no longer only how much we can keep. It is whose people we are becoming through what we give.

The first reading makes that question uncomfortable in a different way. God addresses Cyrus, a powerful king, as someone he has called and strengthened for the sake of his people. Yet the reading twice says that Cyrus does not know him. The king has received a gift without recognizing its giver.

Both facts matter. The good accomplished through Cyrus is real. His power can serve Israel's deliverance. But usefulness is not the same thing as holiness. A successful action does not make everything about its agent good. God can give a ruler a part to play without making that ruler our saviour.

And Cyrus's difficulty can become our own. We can have an ability, perform a service, even help somebody considerably, and still treat the good as proof that we need no help ourselves. We can speak about God while quietly giving ourselves all the credit. The gift is being used. The giver is being forgotten.

The Psalm turns our attention outward: all peoples are called to give glory to the Lord who made the heavens. There is room in that praise for every person who serves. There is no place in it for a servant to take the Creator's place. We can receive good things gratefully without asking them to save us. We can do good work without making ourselves its final explanation.

Paul shows us what that looks like in a community. He remembers the Thessalonians' faith at work, their love bearing labour, their hope enduring. These are people doing something. Their faith is more than an opinion they hold. Their love costs them effort. Their hope continues when continuing is difficult.

And Paul thanks God for them. He does not have to choose between recognizing their work and recognizing God's gift. Thanksgiving joins the two. Grace has borne fruit in actual people whose actual choices matter. God is at work in a love that has become willing to labour.

This brings us to the opening prayer of this Mass. Its concern is sincere service and the disposition of our will. We ask God for help in becoming willing. That prayer reaches a place where simply being told our duty is not enough. We may know what we owe and still resent it. We may do the right thing while arranging to receive the credit. We may offer help while trying to keep control of the person we help.

The Gospel's questioners knew how to say something true. Their intention still needed conversion. Our prayer brings that need before God. His help can reach the choosing itself. We can ask for a heart that serves sincerely, and then take the next faithful step. We need not wait until goodness has become effortless. Paul's love has labour in it.

So take one responsibility that is already yours this week. Perhaps there is a promise you have kept postponing, or a person whose needs have begun to irritate you. Ask what is actually owed. Then ask what you are trying to secure for yourself while doing it. Praise? Control? An excuse to feel superior? Let that question become an occasion to ask God's help, and carry out the good that is yours to do.

This is also why we come to the altar as people who still need to receive. The prayer over the offerings joins our service to our need for purification. Having something to give does not mean that the giver is finished being changed. The gift we bring cannot buy an exemption from mercy. Our offering becomes a place to acknowledge how deeply we need it.

Christ's own gift comes before our response. One of this Sunday's Communion antiphons recalls his giving his life as ransom. In the Gospel setting of that saying, the disciples are competing for places of honour. Jesus directs them to his saving self-gift. He does more than show them an admirable way to behave. He gives the life by which they can be freed and learn to serve.

We do not redeem ourselves by becoming more impressive servants. We receive the Saviour, and our service begins to take its shape from him. At this Eucharist, the whole life asked of us meets the life Christ gives for us. The one who calls for our return has first made return possible.

And our return is still a journey. The prayer after Communion directs present help toward eternal gifts. Much may remain unfinished when we leave: the difficult duty, the mixed motive, the love that needs patience. Hope can endure there. Its fulfilment rests with God, beyond what our own strength or any earthly power can accomplish.

Return to that small coin. Jesus lets it remain a coin, and lets the debt remain a debt. But he does not let the person holding it disappear behind the transaction. That person belongs to God. So do the gifts that make service possible. So does the hope that sustains it. Render the duty that is truly yours. Receive the mercy you still need. Let the life you have received become a life returned in love.
